\section{Introduction}

Optimal-auction theory characterizes the seller's revenue-maximizing reserve
price as a functional of the latent distribution of bidder valuations
\citep{myerson1981}. In empirical settings, however, researchers and sellers
often observe only a transaction price, a winning bid, or a small number of
order statistics rather than the complete vector of private values. This paper
asks whether a flexible generative model can recover the economically relevant
features of the latent distribution while still permitting valid inference on
the reserve-price parameter.

This question is difficult for a reason that is easy to hide with global fit
statistics. If only the highest of $N$ valuations is observed, then the
probability of seeing an auction maximum below the reserve is $F_0(r_0)^N$.
The reserve can therefore depend on a region with very few observations even
when the observable maximum-bid distribution is well approximated overall.
Quantile or Wasserstein accuracy averaged over the support need not control the
local CDF and density errors entering the Myerson first-order condition.

The proposed framework treats the valuation distribution as an
infinite-dimensional nuisance parameter and the reserve price as the
low-dimensional parameter of interest. A structural WGAN-GP reproduces the
auction observation process: it generates independent bidder values, applies
the relevant order-statistic map, and matches the resulting observable
distribution to the data. Cross-fitting separates nuisance estimation from
score evaluation, while a Neyman-orthogonal correction reduces the first-order
sensitivity of the reserve-price estimator to errors in the learned
distribution.

The paper makes three contributions. First, it connects adversarial structural
estimation to the order-statistics approach in empirical auctions. The
generator respects the independent-private-values structure before applying
the same maximum operator observed in the data. Second, it derives an
observation-level orthogonal score for a smoothed reserve and gives
conditions for cross-fitted asymptotic linearity and valid fixed-bandwidth
inference. A covariance-aware Richardson combination provides a separate route
from that regularized target to the exact reserve. Third, it uses nested Monte
Carlo ablations to distinguish failures of global WGAN learning, local nuisance
rates, smoothing bias, and empirical root selection rather than attributing all
errors to the generative model.

The results sharpen the role of WGAN-GP in this problem. It reproduces the
observable maximum distribution reasonably well but misses local tail moments
that matter for the reserve. Replacing only those score nuisances with
training-fold empirical moments restores near-oracle coverage for the
fixed-bandwidth target. For the exact reserve, a stable $n^{-0.2}$ bandwidth
path combined with Richardson extrapolation reduces bias, and a
three-bandwidth residual-bias envelope attains 96.5 percent coverage in the
reported design. More aggressive shrinkage creates multiple roots; taking the
union over all root chains increases interval length without improving
coverage.

These findings support a hybrid estimator rather than a WGAN-only inference
procedure. The WGAN supplies a coherent global latent distribution for
structural counterfactuals, while a dedicated local learner supplies the score
nuisances used for inference. The paper's formal coverage result concerns the
fixed-bandwidth regularized reserve. Exact-reserve intervals are reported as a
bias-aware prototype whose uniform validity requires additional conditions on
the shrinking-bandwidth root process.
